% P21.3 residual-Picard mixed-tower consequences, 11 September 2026.
\begin{theorem}[Residual-Picard descent through mixed locally trivial towers]
\label{thm:p213-residual-picard-mixed-tower}
Let \(K,C,W\) be as in
\Cref{prop:p213-residual-picard-smooth-curve}, assume
\eqref{eq:p213-residual-picard-vanishing}, and put \(B_0=W\).  Let
\(\ell\geq1\), and suppose that
\[
 B_\ell\xrightarrow{\pi_\ell}B_{\ell-1}
 \longrightarrow\cdots\longrightarrow B_1
 \xrightarrow{\pi_1}B_0
\]
is a finite tower in which every \(B_j\) is a positive-dimensional good,
topologically separated, pure strict \(K\)-analytic space.  For every
\(1\leq j\leq\ell\), assume that exactly one of the following descriptions
is specified.
\begin{enumerate}
\item The morphism \(\pi_j\) is a locally trivial standard strict
open-polydisc bundle as in
\Cref{prop:p062-open-polydisc-bundle-descent}; in particular, its positive
fibre dimension and all chart radii satisfy the conditions there.
\item The base \(B_{j-1}\) is locally holomorphically separated and
paracompact, and \(\pi_j\) is a locally trivial analytic projective-space
bundle as in \Cref{prop:p062-projective-bundle-descent}, with fibre
\((\mathbf P_K^{q_j})^{\mathrm{an}}\) for some \(q_j\geq1\).
\item The base $B_{j-1}$ is locally holomorphically separated and
paracompact, and $\pi_j$ is a locally trivial smooth proper split-toric
bundle as in
\Cref{prop:p062-smooth-proper-split-toric-bundle-descent}, with fibre
$X_{\Sigma_j}^{\mathrm{an}}$ for a smooth proper split toric $K$-variety
$X_{\Sigma_j}$ of dimension $q_j\geq1$.
\item The base \(B_{j-1}\) is locally holomorphically separated and
paracompact, and \(\pi_j\) is a locally trivial proper smooth curve bundle
as in \Cref{prop:p062-proper-smooth-curve-bundle-descent}, with fibre
\(D_j^{\mathrm{an}}\) for a nonempty geometrically connected smooth proper
algebraic \(K\)-curve \(D_j\).
\end{enumerate}
If \(h\in C^0(B_\ell,\mathbf R)\) satisfies
\(\dprimeddprime[h]=0\), then there is a unique
\(b\in C^0(W,\mathbf R)\) such that, for
\(\Pi=\pi_1\circ\cdots\circ\pi_\ell\),
\[
                  h=\Pi^*b,
                  \qquad \dprimeddprime[b]=0.
\]
Moreover, \(b\) and \(h\) are ordinary CLD-smooth and have zero smooth
Hessian.
\end{theorem}

\begin{proof}
Starting with \(\pi_\ell\) and proceeding down the tower, apply respectively
\Cref{prop:p062-open-polydisc-bundle-descent,prop:p062-projective-bundle-descent,prop:p062-smooth-proper-split-toric-bundle-descent,prop:p062-proper-smooth-curve-bundle-descent}
according to the specified type of each stage.  The hypotheses above retain
the full geometric assumptions of the applicable descent proposition.  At
every step one obtains a unique continuous zero-current function on the
preceding base.  The pullback identities compose and uniqueness propagates,
giving the displayed function \(b\) on \(W\).

The residual-Picard criterion
\Cref{prop:p213-residual-picard-smooth-curve} makes \(b\) ordinary
CLD-smooth with zero smooth Hessian.  Pullback of ordinary CLD forms is
functorial and commutes with both differentials by
\Cref{prop:cld-form-pullback}.  Therefore \(h=\Pi^*b\) has the same two
properties.
\end{proof}

\begin{corollary}[Residual-Picard mixed-tower Bedford--Taylor criterion]
\label{cor:p213-residual-picard-mixed-tower-bt}
In the setting of \Cref{thm:p213-residual-picard-mixed-tower}, assume that
\(B=B_\ell\) is boundaryless of pure dimension \(n\geq2\).  Let \(T\) be a
positive \(\mathrm d'\)- and \(\mathrm d''\)-closed ordinary CLD current of
bidegree \((p,p)\), \(0\leq p\leq n-1\), and let
\(u,v\in C^0(B,\mathbf R)\) satisfy
\[
                         \dprimeddprime[u]=\dprimeddprime[v].
\]
Then
\[
                         \dprimeddprime(uT)=\dprimeddprime(vT).
\]
In particular, this holds for every continuous psh pair in the
continuous-potential predicate of P21.3.
\end{corollary}

\begin{proof}
Put \(h=u-v\).  The tower theorem makes \(h\) ordinary CLD-smooth with
zero smooth Hessian.  Apply \Cref{cor:ch16-smooth-h}.
\end{proof}

\begin{corollary}[Residual-Picard mixed-tower cover criterion]
\label{cor:p213-residual-picard-mixed-tower-cover-bt}
Let \(K\) be algebraically closed, complete, and nontrivially valued, let
\(X\) be a good, topologically separated, pure, boundaryless strict
\(K\)-analytic space of dimension \(n\geq2\), and let \(U\subseteq X\) be
ordinary open.  Suppose that \(U\) has an ordinary open cover by sets \(B\)
which, with their induced analytic structures, satisfy the geometric
hypotheses of \Cref{cor:p213-residual-picard-mixed-tower-bt}.  Then every
positive closed ordinary CLD current \(T\) of bidegree \((p,p)\) on \(U\),
\(0\leq p\leq n-1\), has Bedford--Taylor independence for every continuous
psh pair \(u,v\) on \(U\) with equal Hessian currents.
\end{corollary}

\begin{proof}
Put \(h=u-v\).  On every member \(B\) of the displayed cover, apply
\Cref{thm:p213-residual-picard-mixed-tower} to \(h|_B\).  Thus \(h\) is
ordinary CLD-smooth with zero smooth Hessian locally on \(U\), and hence has
those properties on \(U\).  Apply \Cref{cor:ch16-smooth-h} directly to
\(T,u,v\) on \(U\).
\end{proof}

\begin{corollary}[Residual-Picard mixed-tower basis criterion]
\label{cor:p213-residual-picard-mixed-tower-basis-bt}
Let \(K\) and \(X\) be as in
\Cref{cor:p213-residual-picard-mixed-tower-cover-bt}.  Suppose that the
ordinary topology of \(X\) has a basis of open sets \(B\) such that every
\(B\), with its induced analytic structure, satisfies the geometric
hypotheses of \Cref{cor:p213-residual-picard-mixed-tower-bt}.  Then the
continuous-pair Bedford--Taylor predicate in P21.3 holds on every ordinary
open subset of \(X\) for every positive closed ordinary CLD current.
\end{corollary}

\begin{proof}
For an ordinary open \(U\subseteq X\), the basis members contained in \(U\)
form a cover to which
\Cref{cor:p213-residual-picard-mixed-tower-cover-bt} applies.
\end{proof}
